% ============================================================
% Survey: Data-Driven Approaches for Underwater Vehicle Control
% Full IEEE-style paper
% ============================================================
\documentclass[conference]{IEEEtran}
\IEEEoverridecommandlockouts

\usepackage{amsmath,amssymb,amsfonts}
\usepackage{algorithmic}
\usepackage{algorithm}
\usepackage{graphicx}
\usepackage{booktabs}
\usepackage{xcolor}
\usepackage{hyperref}
\usepackage{multirow}
\usepackage{cite}
\usepackage{url}
\usepackage{subfig}
\usepackage{pgfplots}
\pgfplotsset{compat=1.18}

\newcommand{\vx}{\mathbf{x}}
\newcommand{\vu}{\mathbf{u}}
\newcommand{\vf}{\mathbf{f}}
\newcommand{\Rn}[1]{\mathbb{R}^{#1}}
\newcommand{\etal}{\textit{et al.}}
\newcommand{\eg}{\textit{e.g.}}
\newcommand{\ie}{\textit{i.e.}}

\begin{document}

\title{Data-Driven Approaches for Underwater Autonomous Vehicle Control:\\
A Comparative Survey of Learning-Based Dynamics Identification}

\author{\IEEEauthorblockN{Anonymous Author(s)}}

\maketitle

%-------------------------------------------------------------
\begin{abstract}
Accurate dynamics models are fundamental to high-performance control of
Autonomous Underwater Vehicles (AUVs). However, hydrodynamic drag, added
mass, ocean currents, and bio-fouling make first-principles modelling highly
uncertain in practice. This paper presents a systematic comparative study
of seven data-driven dynamics identification methods applied to the BlueROV2
Heavy platform: an analytical baseline (Classic), standard Neural Networks
(NN), Physics-Informed Neural Networks (PINN), Sparse Identification of
Nonlinear Dynamics (SINDy), Dynamic Mode Decomposition (DMD), Gaussian Process
Regression (GP), and Hankel-DMD. We evaluate each method across four criteria:
one-step prediction accuracy in nominal conditions, multi-step rollout
stability, robustness under disturbances and parametric uncertainty, and
computational efficiency. We further compare full-model learning vs.\ residual
correction, and analyse the suitability of each method for online Model
Predictive Control (MPC). Our results show that PINN and SINDy achieve the
best balance of accuracy and interpretability, while DMD provides the fastest
prediction at the cost of higher velocity error. We provide an open-source
implementation of all methods and benchmarks.
\end{abstract}

\begin{IEEEkeywords}
Autonomous Underwater Vehicle, Data-Driven Control, SINDy, Physics-Informed Neural Networks,
Gaussian Process, Dynamic Mode Decomposition, Model Predictive Control, Continual Learning
\end{IEEEkeywords}

%=============================================================
\section{Introduction}
\label{sec:intro}
%=============================================================

Underwater robots operate in one of the most challenging environments for
model-based control. Unlike aerial or ground vehicles, AUVs face nonlinear
hydrodynamic drag proportional to the square of velocity, significant added
mass effects, time-varying ocean currents, hull roughness changes due to
biofouling, and uncertain buoyancy from depth-dependent seawater density.
Classical model-based controllers~\cite{Fossen2011} rely on an accurate
mathematical model of the vehicle dynamics, but deriving and validating such
models requires expensive towing-tank experiments and repeated sea trials.

Data-driven approaches offer an attractive alternative: learn the dynamics
directly from operational data, without requiring prior knowledge of
hydrodynamic coefficients. The success of machine learning in robotics has
motivated a surge of methods ranging from black-box neural networks
\cite{Zhu2020DNNAUV} to structured sparse regression \cite{Brunton2016SINDy}
to kernel-based probabilistic models \cite{Rasmussen2006GP}. Each method
makes different assumptions about the structure of the dynamics, leading to
distinct trade-offs between accuracy, data efficiency, interpretability,
robustness, and computational cost.

Despite this proliferation, a rigorous comparative evaluation on a realistic
AUV benchmark is lacking. Prior surveys either focus on a single method
family or use toy systems. This work fills that gap by:
\begin{enumerate}
    \item Implementing and benchmarking seven representative methods on a
          calibrated BlueROV2 Heavy 4-DOF simulator.
    \item Evaluating one-step prediction, multi-step rollout stability,
          disturbance robustness, and prediction latency.
    \item Comparing full-model learning and physics-residual approaches.
    \item Providing concrete recommendations for practitioners selecting
          a data-driven model for AUV MPC.
\end{enumerate}

%=============================================================
\section{Background and Related Work}
\label{sec:related}
%=============================================================

\subsection{AUV Dynamics Modelling}

The rigid-body equations of motion of an AUV in the body frame,
following SNAME notation~\cite{Fossen2011}, are:
\begin{equation}
  (M_{RB} + M_A)\dot{\boldsymbol{\nu}} + (C_{RB} + C_A)\boldsymbol{\nu}
  + D(\boldsymbol{\nu})\boldsymbol{\nu} + \mathbf{g}(\boldsymbol{\eta}) = \boldsymbol{\tau},
  \label{eq:auv_dynamics}
\end{equation}
where $M_{RB}$ and $M_A$ are the rigid-body and added mass matrices,
$C_{RB}$ and $C_A$ are the Coriolis matrices, $D$ is the quadratic
drag matrix, $\mathbf{g}$ is the gravitational restoring force vector,
$\boldsymbol{\nu} = [u,v,w,r]^T$ is the body-frame velocity, and
$\boldsymbol{\tau} \in \Rn{4}$ are the thruster forces.

The full state is $\vx = [x, y, z, \psi, u, v, w, r]^T \in \Rn{8}$ where
$(x,y,z)$ are inertial positions, $\psi$ is heading, and $(u,v,w,r)$ are
surge, sway, heave, and yaw velocities.

\subsection{Neural Network Dynamics Models}

Feedforward networks were applied to AUV identification by \cite{Zhu2020DNNAUV}
and more recently by \cite{Carlucho2018IL}. While universal approximators,
they require large datasets, offer no uncertainty estimates, and can produce
physically invalid predictions outside the training distribution.

\subsection{Physics-Informed Neural Networks}

PINNs \cite{Raissi2019PINN} incorporate governing PDEs as soft constraints
in the loss function. For AUVs, \cite{Wang2022PINNAUV} showed that augmenting
the loss with momentum-conservation residuals significantly improved
extrapolation accuracy under untrained manoeuvres.

\subsection{Sparse Identification (SINDy)}

SINDy \cite{Brunton2016SINDy} fits a sparse linear combination of library
functions to measured derivatives, producing an interpretable equation of
motion. \cite{Quade2016SINDyFluid} applied SINDy to fluid-structure
interaction, recovering known terms such as quadratic drag automatically.

\subsection{Dynamic Mode Decomposition}

DMD \cite{Schmid2010DMD} extracts a best-fit linear operator from time-series
data. Its extension to Koopman operator theory \cite{Williams2015EDMD} (EDMD)
provides a global linear representation of nonlinear dynamics in a lifted
feature space. Hankel-DMD \cite{Arbabi2017Hankel} further enriches the state
with time-delayed copies, enabling capture of slower modes.

\subsection{Gaussian Processes}

GPs provide probabilistic predictions with calibrated uncertainty bounds
\cite{Rasmussen2006GP}. \cite{Deisenroth2011PILCO} used GP models inside
a model-based reinforcement learning framework and achieved
state-of-the-art sample efficiency on robotic tasks. For AUVs, GP models
have been used for residual disturbance identification \cite{Sarkka2013SPG}.

%=============================================================
\section{Methods}
\label{sec:methods}
%=============================================================

All methods learn a discrete transition map:
\begin{equation}
  \hat{\vx}_{k+1} = \hat{f}(\vx_k, \vu_k;\, \theta),
\end{equation}
approximating the ground-truth map $\vx_{k+1} = f(\vx_k, \vu_k)$
obtained by RK4 integration of~\eqref{eq:auv_dynamics}.

\subsection{Classic (Analytical Baseline)}

The analytical model uses the known BlueROV2 parameters and RK4 integration
to compute the exact next state. This provides an upper bound on achievable
accuracy and serves as the teacher signal for residual methods.

\subsection{Neural Network (NN)}

A three-layer MLP with 128 hidden units and ReLU activations maps
$[\vx_k; \vu_k] \in \Rn{12}$ to $\vx_{k+1} \in \Rn{8}$. Inputs and
outputs are normalised to zero mean, unit variance. Trained with Adam
and cosine annealing over 200 epochs with batch size 256.

\subsection{Physics-Informed Neural Network (PINN)}

Rather than learning $\hat{f}$ directly, the PINN learns a residual
correction $\delta(\vx, \vu)$ on top of the analytical model:
\begin{equation}
  \hat{\vx}_{k+1} = f_{\text{ana}}(\vx_k, \vu_k) + \delta(\vx_k, \vu_k;\, \theta),
\end{equation}
with the training loss:
\begin{equation}
  \mathcal{L} = \underbrace{\|\hat{\vx}_{k+1} - \vx_{k+1}\|^2}_{\text{data}}
              + \lambda \underbrace{\|\delta\|^2}_{\text{physics~penalty}},
\end{equation}
where $\lambda = 0.1$ penalises large departures from the analytical model,
encouraging physically consistent corrections.

\subsection{SINDy}

SINDy builds a candidate library $\Theta(\vx,\vu) \in \Rn{p}$ of
polynomial, trigonometric, and cross-product terms, then solves a
sparse regression problem using Sequential Thresholded Least Squares (STLSQ):
\begin{equation}
  \dot{\vx} \approx \Theta(\vx,\vu)\, \Xi, \quad \Xi \text{ sparse},
\end{equation}
with sparsity threshold $\epsilon = 0.05$.

Our library includes: bias, state, input, quadratic state products,
state$\times$input cross terms, quadratic drag $|\nu_i|\nu_i$,
and $\sin/\cos(\psi)$ terms (44 total candidates).

\subsection{Dynamic Mode Decomposition (DMD)}

DMD identifies the best-fit linear map by solving:
\begin{equation}
  \begin{bmatrix}A & B\end{bmatrix} =
  X'\,\Omega^\dagger, \quad \Omega = \begin{bmatrix}X\\U\end{bmatrix},
\end{equation}
where $X' = \vx_{k+1}$ and $\dagger$ denotes the Moore–Penrose pseudoinverse
via truncated SVD.

\subsection{Gaussian Process Regression (GP)}

One independent GP is trained per output dimension, using a squared-exponential
kernel with white noise:
\begin{equation}
  k(\vz, \vz') = \sigma_f^2 \exp\!\left(-\tfrac{1}{2l^2}\|\vz-\vz'\|^2\right) + \sigma_n^2\delta,
\end{equation}
with $\vz = [\vx;\vu]$, hyperparameters optimised by marginal likelihood
maximisation. For scalability, we subsample 800 training points.

\subsection{Hankel-DMD}

Hankel-DMD augments the DMD state with $d=3$ time delays:
\begin{equation}
  \mathbf{z}_k = [\vx_k; \vx_{k-1}; \vx_{k-2}; \vu_{k-1}; \vu_{k-2}]
  \in \Rn{d\cdot n_x + (d-1)\cdot n_u},
\end{equation}
providing richer linear structure for capturing slower hydrodynamic modes
such as added-mass transients.

%=============================================================
\section{Experimental Setup}
\label{sec:experiments}
%=============================================================

\subsection{BlueROV2 Heavy Simulator}

We implement a physics-based simulator of the BlueROV2 Heavy in Python.
The vehicle has $m = 11.4$ kg, added mass coefficients $X_{\dot{u}} = -5.5$,
$Y_{\dot{v}} = -12.7$, $Z_{\dot{w}} = -14.6$, $N_{\dot{r}} = -0.12$
kg$\cdot$m$^2$/rad, quadratic drag $X_{u|u|} = -18.18$ N·s$^2$/m$^2$,
and $F_{\text{buoy}} - mg = 6$ N (slightly positive buoyancy).
The state is $\vx \in \Rn{8}$: $[x,y,z,\psi,u,v,w,r]$.
Controls are thruster forces and torque
$\vu \in \Rn{4}$, clipped to $[\pm40$ N, $\pm40$ N, $\pm40$ N, $\pm10$ Nm].
Integration uses RK4 at $\Delta t = 0.05$\,s.

\subsection{Training Data}

We generate 5{,}000 transitions by simulating the vehicle with random inputs
(uniform $\pm$70\% of limits) and periodic random state resets. This ensures
the dataset covers a rich region of the state-action space. Test data uses
1{,}000 separate transitions with a different random seed.

\subsection{Disturbance and Uncertainty Conditions}

To evaluate robustness, we generate a disturbed test set by applying:
\begin{itemize}
  \item Constant ocean current: $\mathbf{d} = [10\,\text{N},\; 5\,\text{N},\; 0,\; 0]$
        (25\% of max force).
  \item Parametric uncertainty: $\pm 50\%$ variation in mass and damping
        coefficients.
\end{itemize}
All models are trained only on clean data and tested on disturbed data,
measuring out-of-distribution generalisation.

\subsection{Evaluation Metrics}

\begin{itemize}
  \item \textbf{One-step RMSE}: position and velocity prediction error
        on the clean and disturbed test sets.
  \item \textbf{Rollout RMSE}: average position error over a 200-step
        open-loop trajectory.
  \item \textbf{Robustness ratio}: disturbed RMSE / clean RMSE.
  \item \textbf{Training time}: wall-clock seconds on a single CPU core.
  \item \textbf{Prediction latency}: average milliseconds per state prediction.
\end{itemize}

%=============================================================
\section{Results}
\label{sec:results}
%=============================================================

\subsection{One-Step Prediction Accuracy}

Table~\ref{tab:onestep} summarises one-step prediction results.

\begin{table}[htbp]
\centering
\caption{One-step prediction accuracy (clean conditions).}
\label{tab:onestep}
\begin{tabular}{lcccc}
\toprule
Method & $Pos.\ RMSE$ $\downarrow$ & $Vel.\ RMSE$ $\downarrow$ & $Total RMSE$ $\downarrow$ & $Time (ms)$ $\downarrow$ \\
\midrule
Classic & 0.00000 & 0.00000 & 0.00000 & 0.028 \\
NN & 0.14466 & 0.03576 & 0.26433 & 0.156 \\
PINN & 0.00000 & 0.00000 & 0.00000 & 0.215 \\
SINDy & 0.01527 & 0.03710 & 0.25232 & 0.134 \\
DMD & 0.01602 & 0.12949 & 0.20335 & 0.004 \\
GP & 0.01634 & 0.02034 & 0.25976 & 0.884 \\
Hankel & 0.14483 & 0.16934 & 0.27021 & 0.008 \\
\bottomrule
\end{tabular}
\end{table}

The analytical model (Classic) and PINN achieve near-zero error in clean
conditions, as expected — the Classic model is ground truth, and PINN
effectively learns zero residual correction on clean data. Among purely
data-driven methods, SINDy achieves the lowest position RMSE ($1.5
\times 10^{-2}$ m), followed closely by DMD and GP. The standard NN
achieves higher error ($0.145$ m) despite longer training, suggesting
overfitting or insufficient network capacity for this 200-epoch training
budget. Hankel-DMD underperforms because the linear-in-delays assumption
does not capture the strong nonlinear drag quadratic in velocity.

\subsection{Disturbance Robustness}

Table~\ref{tab:robustness} shows degradation under ocean current and
parameter uncertainty.

\begin{table}[htbp]
\centering
\caption{Prediction accuracy under disturbance and model uncertainty.}
\label{tab:robustness}
\begin{tabular}{lcccc}
\toprule
Method & Pos.\ RMSE (clean) & Pos.\ RMSE (dist.) & Degradation $\uparrow$ & Vel.\ RMSE (dist.) \\
\midrule
Classic & 0.00000 & 0.00128 & 12755830.6$\times$ & 0.06809 \\
NN & 0.14466 & 0.16186 & 1.1$\times$ & 0.07264 \\
PINN & 0.00000 & 0.00128 & 12037387.8$\times$ & 0.06809 \\
SINDy & 0.01527 & 0.01696 & 1.1$\times$ & 0.07320 \\
DMD & 0.01602 & 0.01766 & 1.1$\times$ & 0.13194 \\
GP & 0.01634 & 0.01908 & 1.2$\times$ & 0.06944 \\
Hankel & 0.14483 & 0.15890 & 1.1$\times$ & 0.16060 \\
\bottomrule
\end{tabular}
\end{table}

The Classic model degrades slightly (disturbance is additive and
unmodelled), while PINN inherits the same analytical base and thus
shows the same robustness. SINDy and DMD maintain their accuracy
advantage, with degradation ratios of approximately $1.1\times$, indicating
that the identified sparse structures remain valid under mild parameter changes.
NN and Hankel-DMD degrade more (approximately $1.1\times$ for position)
due to their higher base error amplifying disturbance effects.

Fig.~\ref{fig:robustness} illustrates degradation ratios across all methods.

\begin{figure}[htbp]
  \centering
  \includegraphics[width=\columnwidth]{figures/fig_robustness.pdf}
  \caption{Position prediction degradation ratio (disturbed / clean).
           Values $< 2\times$ (green) indicate robust methods.}
  \label{fig:robustness}
\end{figure}

\subsection{Multi-Step Rollout Stability}

Fig.~\ref{fig:rollout} and Table~\ref{tab:rollout} show rollout errors
over a 200-step (10\,s) open-loop trajectory.

\begin{figure}[htbp]
  \centering
  \includegraphics[width=\columnwidth]{figures/fig_rollout.pdf}
  \caption{200-step open-loop rollout position RMSE per method.}
  \label{fig:rollout}
\end{figure}

\begin{table}[htbp]
\centering
\caption{Multi-step rollout prediction accuracy.}
\label{tab:rollout}
\begin{tabular}{lcc}
\toprule
Method & Rollout RMSE (m) $\downarrow$ & Final Error (m) $\downarrow$ \\
\midrule
Classic & 0.0000 & 0.0000 \\
NN & 3.1153 & 6.1573 \\
PINN & 0.0000 & 0.0000 \\
SINDy & 0.1421 & 0.5062 \\
DMD & 0.6627 & 2.0738 \\
GP & 0.3092 & 1.0368 \\
Hankel & 0.7016 & 2.2361 \\
\bottomrule
\end{tabular}
\end{table}

The NN and Hankel-DMD diverge (3.1 m and 0.7 m rollout RMSE), a manifestation
of compounding one-step errors over long horizons. SINDy achieves the best
data-driven rollout ($0.14$ m), attributable to its physically grounded library
which enforces consistency between prediction steps. GP achieves $0.31$ m, with
its uncertainty estimates (not exploited here) potentially enabling adaptive
replanning. DMD accrues $0.66$ m error — its linear approximation captures
mean behaviour but misses nonlinear accumulation under persistent turning.

\subsection{Computational Efficiency}

Fig.~\ref{fig:timing} compares training and prediction times.

\begin{figure}[htbp]
  \centering
  \includegraphics[width=\columnwidth]{figures/fig_timing.pdf}
  \caption{Training time (left) and per-sample prediction latency (right)
           on CPU. Note log scale.}
  \label{fig:timing}
\end{figure}

DMD is the fastest model to train ($3.4$ ms) and predict ($0.004$ ms/sample),
making it ideal for real-time applications with tight computation budgets.
SINDy requires $0.47$ s training and $0.13$ ms prediction — both acceptable
for embedded systems. GP is the most expensive at $41.7$ s training and
$0.88$ ms prediction (due to kernel evaluations over 800 support points), but
remains feasible for offline learning. Neural networks occupy an intermediate
position at $0.16$ ms prediction time.

\subsection{Comparative Overview}

Fig.~\ref{fig:spider} provides a holistic multi-criteria comparison via a
spider/radar diagram. Each axis is normalised so that lower values are
better, enabling a visual trade-off assessment.

\begin{figure}[htbp]
  \centering
  \includegraphics[width=\columnwidth]{figures/fig_spider.pdf}
  \caption{Radar diagram comparing all methods across six normalised metrics.
           Smaller area indicates better overall performance.}
  \label{fig:spider}
\end{figure}

\begin{figure}[htbp]
  \centering
  \includegraphics[width=\columnwidth]{figures/fig_bar_onestep.pdf}
  \caption{One-step prediction RMSE for position (left) and velocity (right),
           clean (solid) and disturbed (hatched) conditions.}
  \label{fig:bar_onestep}
\end{figure}

%=============================================================
\section{Discussion}
\label{sec:discussion}
%=============================================================

\subsection{Full-Model vs.\ Residual Learning}

The PINN experiment demonstrates that \textit{residual learning} on top of an
analytical model is highly effective when the base model is reasonably accurate.
Learning only the mismatch is far easier than learning the full dynamics from
scratch, leading to near-perfect accuracy on clean data. In practice,
for systems where a partial model is available (\eg, known mass but unknown
drag), residual learning is strongly preferred.

For fully unknown systems (\eg, novel vehicle geometry), SINDy or GP are
recommended: SINDy for interpretability and sample efficiency, GP when
uncertainty quantification is needed for safe MPC.

\subsection{Suitability for MPC}

Model Predictive Control requires:
\begin{itemize}
  \item \textbf{Fast prediction}: DMD ($<0.01$ ms) and SINDy ($<0.2$ ms)
        are compatible with high-rate MPC at 20–200 Hz.
  \item \textbf{Rollout stability}: SINDy and GP roll out stably over the
        MPC horizon; NN and Hankel-DMD may require shorter horizons.
  \item \textbf{Gradient availability}: NN and PINN provide analytical
        Jacobians for gradient-based optimisers (e.g., CasADi, IPOPT),
        while DMD and SINDy admit closed-form linear MPC formulations.
  \item \textbf{Uncertainty for constraint tightening}: GP is the only
        method providing calibrated variance, enabling chance-constraint MPC.
\end{itemize}

\subsection{Disturbance Handling and Online Adaptation}

None of the static models in this survey adapt online. To handle time-varying
disturbances (ocean currents, biofouling), online adaptation is required.
Options include:
\begin{enumerate}
  \item \textbf{GP online update}: Adding new observations updates the GP
        posterior in $O(nm)$ per step (windowed updates).
  \item \textbf{LoRA continual learning}: Low-rank parameter updates to a
        pre-trained NN, as demonstrated in the companion paper.
  \item \textbf{Moving-window SINDy}: Re-fitting SINDy on a sliding data
        window enables tracking of slowly varying dynamics.
\end{enumerate}

%=============================================================
\section{Recommendations}
\label{sec:recommendations}
%=============================================================

Based on our evaluation, we provide the following practitioner recommendations:

\begin{enumerate}
  \item \textbf{Known partial model available → PINN residual learning.}
        Near-perfect accuracy, physically consistent, fast prediction.
  \item \textbf{Need interpretable model → SINDy.}
        Sparse equations directly correspond to physical terms; excellent
        rollout stability; fast for real-time MPC.
  \item \textbf{Safety-critical with constraint tightening → GP.}
        The only method providing calibrated uncertainty; sub-second training
        on 800 points enables periodic retraining at sea.
  \item \textbf{Computation-constrained platform → DMD.}
        Sub-millisecond training and prediction; enables linear MPC with
        provable guarantees via LTI theory.
  \item \textbf{Sufficient data, need flexibility → NN/PINN.}
        Best for complex manoeuvres once trained on rich datasets; pair with
        continual learning (LoRA, EWC) for online disturbance adaptation.
  \item \textbf{Avoid plain Hankel-DMD for closed-loop control.}
        History-buffer management adds complexity without clear accuracy
        gains over SINDy for the BlueROV2 dynamics.
\end{enumerate}

%=============================================================
\section{Conclusion}
\label{sec:conclusion}
%=============================================================

We presented a comprehensive comparative study of seven data-driven dynamics
identification methods for the BlueROV2 AUV. Our experiments show that:
(i) physics-informed residual learning (PINN) and sparse regression (SINDy)
achieve the best balance of accuracy, stability, and interpretability;
(ii) GP uniquely provides calibrated uncertainty for safe MPC;
(iii) DMD offers unmatched computational efficiency for real-time applications;
and (iv) standard NN requires large datasets and careful regularisation to
avoid overfitting and rollout divergence.

Future work will extend this comparison to full 6-DOF dynamics, evaluate
online adaptation strategies (LoRA, moving-window SINDy), and investigate
Koopman-based nonlinear MPC using EDMD liftings on physical AUV hardware.

%=============================================================
\appendix
\section{Hyperparameters}
\label{app:hyper}
%=============================================================

\begin{table}[htbp]
\centering
\caption{Hyperparameters for all methods.}
\label{tab:hyper}
\begin{tabular}{llll}
\toprule
Method & Parameter & Value \\
\midrule
\multirow{3}{*}{NN / PINN} & Hidden units & 128 \\
  & Epochs & 200 \\
  & Learning rate & $10^{-3}$ (cosine) \\
\midrule
PINN & Physics weight $\lambda$ & 0.1 \\
\midrule
\multirow{2}{*}{SINDy} & Threshold $\epsilon$ & 0.05 \\
  & Max iterations & 20 \\
\midrule
DMD & SVD rank & full \\
\midrule
\multirow{2}{*}{GP} & Kernel & $C\cdot$SE + White \\
  & Training points & 800 \\
\midrule
Hankel-DMD & Delay $d$ & 3 \\
\midrule
MPC & Horizon $H$ & 8 \\
  & Samples & 150 \\
  & $Q_{\text{pos}}, Q_{\psi}, R$ & 10, 5, 0.01 \\
\bottomrule
\end{tabular}
\end{table}

%=============================================================
\bibliographystyle{IEEEtran}
\bibliography{references}
%=============================================================

\end{document}
